\documentclass[10pt,a4paper]{amsart}
\usepackage{amsmath,amssymb,amsthm,graphicx}
\usepackage{verbatim}
\newenvironment{reference}{\comment}{\endcomment}
%\newenvironment{reference}{}{}
\newenvironment{slogan}{\comment}{\endcomment}
\newenvironment{history}{\comment}{\endcomment}

% For commutative diagrams we use Xy-pic
\usepackage[all]{xy}

% We use 2cell for 2-commutative diagrams.
\xyoption{2cell}
% Theorem environments.
%
\theoremstyle{plain}
\newtheorem{theorem}[subsection]{Theorem}
\newtheorem{proposition}[subsection]{Proposition}
\newtheorem{lemma}[subsection]{Lemma}

\theoremstyle{definition}
\newtheorem{definition}[subsection]{Definition}
\newtheorem{example}[subsection]{Example}
\newtheorem{exercise}[subsection]{Exercise}
\newtheorem{situation}[subsection]{Situation}

\theoremstyle{remark}
\newtheorem{remark}[subsection]{Remark}
\newtheorem{remarks}[subsection]{Remarks}

\numberwithin{equation}{subsection}

% Macros
%
\def\lim{\mathop{\mathrm{lim}}\nolimits}
\def\colim{\mathop{\mathrm{colim}}\nolimits}
\def\Spec{\mathop{\mathrm{Spec}}}
\def\Hom{\mathop{\mathrm{Hom}}\nolimits}
\def\Ext{\mathop{\mathrm{Ext}}\nolimits}
\def\SheafHom{\mathop{\mathcal{H}\!\mathit{om}}\nolimits}
\def\SheafExt{\mathop{\mathcal{E}\!\mathit{xt}}\nolimits}
\def\Sch{\mathit{Sch}}
\def\Mor{\mathop{\mathrm{Mor}}\nolimits}
\def\Ob{\mathop{\mathrm{Ob}}\nolimits}
\def\Sh{\mathop{\mathit{Sh}}\nolimits}
\def\NL{\mathop{N\!L}\nolimits}
\def\CH{\mathop{\mathrm{CH}}\nolimits}
\def\proetale{{pro\text{-}\acute{e}tale}}
\def\etale{{\acute{e}tale}}
\def\QCoh{\mathit{QCoh}}
\def\Ker{\mathop{\mathrm{Ker}}}
\def\Im{\mathop{\mathrm{Im}}}
\def\Coker{\mathop{\mathrm{Coker}}}
\def\Coim{\mathop{\mathrm{Coim}}}

% Boxtimes
%
\DeclareMathSymbol{\boxtimes}{\mathbin}{AMSa}{"02}

%
% Macros for moduli stacks/spaces
%
\def\QCohstack{\mathcal{QC}\!\mathit{oh}}
\def\Cohstack{\mathcal{C}\!\mathit{oh}}
\def\Spacesstack{\mathcal{S}\!\mathit{paces}}
\def\Quotfunctor{\mathrm{Quot}}
\def\Hilbfunctor{\mathrm{Hilb}}
\def\Curvesstack{\mathcal{C}\!\mathit{urves}}
\def\Polarizedstack{\mathcal{P}\!\mathit{olarized}}
\def\Complexesstack{\mathcal{C}\!\mathit{omplexes}}
% \Pic is the operator that assigns to X its picard group, usage \Pic(X)
% \Picardstack_{X/B} denotes the Picard stack of X over B
% \Picardfunctor_{X/B} denotes the Picard functor of X over B
\def\Pic{\mathop{\mathrm{Pic}}\nolimits}
\def\Picardstack{\mathcal{P}\!\mathit{ic}}
\def\Picardfunctor{\mathrm{Pic}}
\def\Deformationcategory{\mathcal{D}\!\mathit{ef}}
\usepackage{cleveref}
\crefname{originalequation}{original equation}{original equations}
\Crefname{originalequation}{Original equation}{Original equations}
\makeatletter
\newlabel{equation-C-prime}{{0AL4}{source}{}{sourcepointer-1}{}}
\newlabel{equation-C-prime@cref}{{[originalequation][0][]0AL4}{[1][0][]source}{}{}{}}
\makeatother
\crefname{originalreference}{original reference}{original references}
\Crefname{originalreference}{Original reference}{Original references}
\makeatletter
\newlabel{remark-improve-homomorphism}{{0GAP}{source}{}{sourcepointer-2}{}}
\newlabel{remark-improve-homomorphism@cref}{{[originalreference][0][]0GAP}{[1][0][]source}{}{}{}}
\makeatother
\title{Lifting algebra maps through powers of a principal ideal with a nonzerodivisor}
\author{The Stacks Project Authors}
\date{stacks-project:restricted.tex:local-snapshot-20261012; GNU FDL 1.2 or later}
\begin{document}
\maketitle
\noindent\textit{Editorial note.} All delivered mathematical text is the author's own, copied from the pinned Stacks Project chapter file restricted.tex; no statement or proof step was rewritten, reordered or removed.\par\medskip
\noindent\textit{Editorial note.} The counted proof is the author's complete argument: the proof identifies the coefficient shift through the nonzerodivisor property, pushes the obstruction past a\textasciicircum{}\{n-c\} using the annihilation hypothesis, observes the exponent grows, iterates, and takes the limit in the complete ring\par\medskip
\noindent\textit{Editorial note.} The label lemma-get-morphism-nonzerodivisor names the counted statement, and the proof environment immediately following it in the source is delivered verbatim as its proof.\par\medskip
\noindent\textit{Editorial note.} The counted statement is carried by the author's own lemma environment in the chapter source.\par\medskip
\noindent\textit{Allowed standard prerequisites (editorial).}\begin{itemize}
\item the obstruction class o(psi\_n) in Ext\textasciicircum{}1\_B(NL\textasciicircum{}wedge\_\{B/A\}, a\textasciicircum{}nC/a\textasciicircum{}\{2n\}C) attached to a map psi\_n : B -> C/a\textasciicircum{}nC, from the same chapter
\item multiplication by a\textasciicircum{}c on a\textasciicircum{}nC/a\textasciicircum{}\{2n\}C identifies with the projection to a\textasciicircum{}\{n-c\}C/a\textasciicircum{}\{2n-c\}C when a is a nonzerodivisor on C
\item the conversion of vanishing obstruction classes into lifted maps of Remark improve-homomorphism of the same chapter
\item I-adic completeness of C and limits of compatible algebra maps
\end{itemize}\medskip
\ifdefined\hypertarget\hypertarget{sourcepointer-1}{}\fi
\noindent\textit{Source reference.} Original equation 0AL4: sources/stacks-restricted/pinned-original-source.tex, source lines 326--326. Prerequisite pointer invoked at proof line(s) 1244; it names a standard fact of the same chapter that the counted argument uses, and every step of the counted proof built on it is retained verbatim in the delivered body.\par\medskip
\ifdefined\hypertarget\hypertarget{sourcepointer-2}{}\fi
\noindent\textit{Source reference.} Original reference 0GAP: sources/stacks-restricted/pinned-original-source.tex, source lines 1101--1101. Prerequisite pointer invoked at proof line(s) 1259; it names a standard fact of the same chapter that the counted argument uses, and every step of the counted proof built on it is retained verbatim in the delivered body.\par\medskip
\begin{lemma}
\label{lemma-get-morphism-nonzerodivisor}
Let $I = (a)$ be a principal ideal of a Noetherian ring $A$.
Let $B$ be an object of (\ref{equation-C-prime}).
Assume given an integer $c \geq 0$ such that $\Ext^1_B(\NL_{B/A}^\wedge, N)$
is annihilated by $a^c$ for all $B$-modules $N$.
Let $C$ be an $I$-adically complete $A$-algebra such that
$a$ is a nonzerodivisor on $C$. Let $n > 2c$. For any $A$-algebra
map $\psi_n : B \to C/a^nC$ there exists an $A$-algebra
map $\varphi : B \to C$ such that
$\psi_n \bmod a^{n - c}C = \varphi \bmod a^{n - c}C$.
\end{lemma}

\begin{proof}
Consider the obstruction class
$$
o(\psi_n) \in \Ext^1_B(\NL_{B/A}^\wedge, a^nC/a^{2n}C)
$$
of Remark \ref{remark-improve-homomorphism}. Since $a$ is a nonzerodivisor
on $C$ the map $a^c : a^nC/a^{2n}C \to a^nC/a^{2n}C$ is isomorphic to the
map $a^nC/a^{2n}C \to a^{n - c}C/a^{2n - c}C$ in the category of $C$-modules.
Hence by our assumption on $\NL_{B/A}^\wedge$
we conclude that the class $o(\psi_n)$ maps to zero in
$$
\Ext^1_B(\NL_{B/A}^\wedge, a^{n - c}C/a^{2n - c}C)
$$
and a fortiori in
$$
\Ext^1_B(\NL_{B/A}^\wedge, a^{n - c}C/a^{2n - 2c}C)
$$
By the discussion in Remark \ref{remark-improve-homomorphism} we obtain a map
$$
\psi_{2n - 2c} : B \to C/a^{2n - 2c}C
$$
which agrees with $\psi_n$ modulo $a^{n - c}C$.
Observe that $2n - 2c > n$ because $n > 2c$.

\medskip\noindent
We may repeat this procedure. Starting with $n_0 = n$ and
$\psi^0 = \psi_n$ we end up getting a strictly increasing
sequence of integers
$$
n_0 < n_1 < n_2 < \ldots
$$
and $A$-algebra homorphisms $\psi^i : B \to C/a^{n_i}C$
such that $\psi^{i + 1}$ and $\psi^i$ agree modulo $a^{n_i - c}C$.
Since $C$ is $I$-adically complete we can take $\varphi$
to be the limit of the maps
$\psi^i \bmod a^{n_i - c}C : B \to C/a^{n_i - c}C$
and the lemma follows.
\end{proof}

\end{document}
